% GENERATED FILE. Edit canonical lessons, then run npm run generate:books.
% canonical-source-hash: fnv1a64:4547b632392c0bc2
% canonical-lessons: LA-C204-affirmo, LA-C204-doleo, LA-C204-orno, LA-C204-conicio, LA-C204-desino

\chapter{To state firmly, To hurt, To decorate, To throw together, To cease}
\label{ch:la-to-state-firmly-to-hurt-to-decorate-to-throw-together-to-cease}
\hlchaptermodality{204}

\begin{chapteropening}
\textbf{By the end of this chapter:} \emph{I can say to state firmly, to hurt, to decorate, to throw together and to cease.}

Complete the last lesson of chapter 204: I can say to state firmly, to hurt, to decorate, to throw together and to cease.
\end{chapteropening}

\section[affirmō, affirmāre]{affirmō, affirmāre — to state firmly}

\label{lesson:LA-C204-affirmo}



\begin{quote}
Before the new one: say the Latin for to annoy, then the Latin for to forget.
\end{quote}

\subsection*{You'll want to know: affirmō, affirmāre}
\textbf{affirmō, affirmāre} — \textquotedblleft{}to state firmly\textquotedblright{}.

\begin{cousinweb}
\textbf{affirmō} is \emph{ad}, \textquotedblleft{}to\textquotedblright{}, with \emph{firmō}, \textquotedblleft{}to make firm\textquotedblright{} (EtymDB 2.1), from \emph{firmus}, \textquotedblleft{}firm\textquotedblright{}: to make a statement firm.
\end{cousinweb}

\begin{grammarlens}[title={What you've built}]
One more word for this chapter.
\end{grammarlens}

\subsection*{Your turn}
\begin{itemize}
\raggedright
  \item \emph{Say it:} \emph{affirmō, affirmāre}
  \item \emph{Say it:} \emph{affirmō, affirmāre}, once more
  \item \emph{Say it:} say \emph{affirmō, affirmāre}
  \item \emph{Recall:} say the Latin for to annoy, then the Latin for to forget, then say \emph{affirmō, affirmāre} again
\end{itemize}

\begin{tcolorbox}[breakable,colback=teal!4,colframe=teal!35!black,title={Before you move on}]
What does affirmō, affirmāre mean? (\textquotedblleft{}To state firmly\textquotedblright{}.) Say it once more.
\end{tcolorbox}
\section[doleō, dolēre]{doleō, dolēre — to hurt, to grieve}

\label{lesson:LA-C204-doleo}



\begin{quote}
Before the new one: say the Latin for to forget, then the Latin for to state firmly.
\end{quote}

\subsection*{You'll want to know: doleō, dolēre}
\textbf{doleō, dolēre} — \textquotedblleft{}to hurt, to grieve\textquotedblright{}.

\begin{cousinweb}
\textbf{doleō} is inherited, usually traced to Proto-Indo-European \emph{*delh1-}, \textquotedblleft{}to split, to cut\textquotedblright{} (EtymDB 2.1): pain as a cutting.
\end{cousinweb}

\begin{grammarlens}[title={What you've built}]
One more word for this chapter.
\end{grammarlens}

\subsection*{Your turn}
\begin{itemize}
\raggedright
  \item \emph{Say it:} \emph{doleō, dolēre}
  \item \emph{Say it:} \emph{doleō, dolēre}, once more
  \item \emph{Say it:} say \emph{doleō, dolēre}
  \item \emph{Recall:} say the Latin for to forget, then the Latin for to state firmly, then say \emph{doleō, dolēre} again
\end{itemize}

\begin{tcolorbox}[breakable,colback=teal!4,colframe=teal!35!black,title={Before you move on}]
What does doleō, dolēre mean? (\textquotedblleft{}To hurt, to grieve\textquotedblright{}.) Say it once more.
\end{tcolorbox}
\section[ōrnō, ōrnāre]{ōrnō, ōrnāre — to decorate}

\label{lesson:LA-C204-orno}



\begin{quote}
Before the new one: say the Latin for to state firmly, then the Latin for to hurt.
\end{quote}

\subsection*{You'll want to know: ōrnō, ōrnāre}
\textbf{ōrnō, ōrnāre} — \textquotedblleft{}to decorate\textquotedblright{}.

\begin{cousinweb}
\textbf{ōrnō} is often explained as a shortened \emph{*ordināre}, \textquotedblleft{}to put in order\textquotedblright{}, from \emph{ōrdō}, \textquotedblleft{}a row, an order\textquotedblright{}, but this is uncertain.
\end{cousinweb}

\begin{grammarlens}[title={What you've built}]
One more word for this chapter.
\end{grammarlens}

\subsection*{Your turn}
\begin{itemize}
\raggedright
  \item \emph{Say it:} \emph{ōrnō, ōrnāre}
  \item \emph{Say it:} \emph{ōrnō, ōrnāre}, once more
  \item \emph{Say it:} say \emph{ōrnō, ōrnāre}
  \item \emph{Recall:} say the Latin for to state firmly, then the Latin for to hurt, then say \emph{ōrnō, ōrnāre} again
\end{itemize}

\begin{tcolorbox}[breakable,colback=teal!4,colframe=teal!35!black,title={Before you move on}]
What does ōrnō, ōrnāre mean? (\textquotedblleft{}To decorate\textquotedblright{}.) Say it once more.
\end{tcolorbox}
\section[coniciō, conicere]{coniciō, conicere — to throw together; to guess}

\label{lesson:LA-C204-conicio}



\begin{quote}
Before the new one: say the Latin for to hurt, then the Latin for to decorate.
\end{quote}

\subsection*{You'll want to know: coniciō, conicere}
\textbf{coniciō, conicere} — \textquotedblleft{}to throw together; to guess\textquotedblright{}.

\begin{cousinweb}
\textbf{coniciō} is \emph{con-}, \textquotedblleft{}together\textquotedblright{}, with \emph{iaciō}, \textquotedblleft{}to throw\textquotedblright{} (EtymDB 2.1): to throw together, then to put facts together, to guess. English \emph{conjecture} comes from it.
\end{cousinweb}

\begin{grammarlens}[title={What you've built}]
One more word for this chapter.
\end{grammarlens}

\subsection*{Your turn}
\begin{itemize}
\raggedright
  \item \emph{Say it:} \emph{coniciō, conicere}
  \item \emph{Say it:} \emph{coniciō, conicere}, once more
  \item \emph{Say it:} say \emph{coniciō, conicere}
  \item \emph{Recall:} say the Latin for to hurt, then the Latin for to decorate, then say \emph{coniciō, conicere} again
\end{itemize}

\begin{tcolorbox}[breakable,colback=teal!4,colframe=teal!35!black,title={Before you move on}]
What does coniciō, conicere mean? (\textquotedblleft{}To throw together; to guess\textquotedblright{}.) Say it once more.
\end{tcolorbox}
\section[dēsinō, dēsinere]{dēsinō, dēsinere — to cease}

\label{lesson:LA-C204-desino}



\begin{quote}
Before the new one: say the Latin for to decorate, then the Latin for to throw together; to guess.
\end{quote}

\subsection*{You'll want to know: dēsinō, dēsinere}
\textbf{dēsinō, dēsinere} — \textquotedblleft{}to cease\textquotedblright{}.

\begin{cousinweb}
\textbf{dēsinō} is \emph{dē}, \textquotedblleft{}off, away\textquotedblright{}, with \emph{sinō}, \textquotedblleft{}to let, to leave\textquotedblright{}: to leave off. \emph{Sinō} is inherited, but its deeper origin is uncertain: EtymDB 2.1 records two different roots.
\end{cousinweb}

\begin{grammarlens}[title={What you've built}]
One more word for this chapter.
\end{grammarlens}

\subsection*{Your turn}
\begin{itemize}
\raggedright
  \item \emph{Say it:} \emph{dēsinō, dēsinere}
  \item \emph{Say it:} \emph{dēsinō, dēsinere}, once more
  \item \emph{Say it:} say \emph{dēsinō, dēsinere}
  \item \emph{Recall:} say the Latin for to decorate, then the Latin for to throw together; to guess, then say \emph{dēsinō, dēsinere} again
\end{itemize}

\begin{tcolorbox}[breakable,colback=teal!4,colframe=teal!35!black,title={Before you move on}]
What does dēsinō, dēsinere mean? (\textquotedblleft{}To cease\textquotedblright{}.) Say it once more.
\end{tcolorbox}
